% ============================================================
%  Uomini senza petto: il coraggio di un giudizio proprio
%  Parte IV: La bussola del pensiero
% ============================================================

\chapter{Uomini senza petto: il coraggio di un giudizio proprio}
\label{cap:p3c24}

\epigraph{«Per ogni allievo che deve essere protetto da un debole eccesso di sensibilità ce ne sono tre che devono essere svegliati dal sonno della fredda volgarità.»}{C. S. Lewis\index[main]{Lewis, C. S.}, \textit{L'abolizione dell'uomo} (1943)}

Il capitolo precedente ha distribuito la responsabilità tra chi progetta, chi regola e chi usa. Ma la parte di chi usa presuppone qualcosa che abbiamo dato per scontato, cioè una persona disposta a formulare un giudizio e a risponderne. Ed è una disposizione che si cede con facilità. La si delega alla macchina, che risponde con tono sicuro (il bias di automazione del cap.~\ref{cap:p3c11}). La si delega al gruppo, perché se tutti la pensano così qualcosa sapranno (il bias del gregge del cap.~\ref{cap:p1c9}). Oppure la si abbandona del tutto, con un'alzata di spalle. Le macchine che orienteranno il futuro imparano anche da queste rinunce, perché un giudizio che nessuno ha voluto firmare diventa comunque un dato.

Questo capitolo propone una tesi, che va presentata come tale. Per sostenere un giudizio proprio non basta la testa. Serve anche quella parte di noi che gli antichi chiamavano \textit{thymos}, e che una lunga tradizione educativa ha imparato a trattare come rumore.

\section{Il cinismo come bias}

L'alzata di spalle ha una forma riconoscibile. È la convinzione che nulla valga davvero la pena e che entusiasmarsi sia da ingenui. È un cinismo comodo, che non richiede fatica e si presenta come lucidità.

Alla prova dei dati, però, assomiglia più a un bias che a una forma di intelligenza. In una serie di studi su grandi campioni in molti Paesi, chi manifestava atteggiamenti più cinici (l'idea che gli altri agiscano quasi sempre per egoismo e secondi fini) otteneva in media punteggi \textit{più bassi}, non più alti, nei test di abilità cognitive. Sono associazioni modeste e correlazionali, non prove che il cinismo renda meno intelligenti. Eppure, quando si chiedeva di immaginare una persona cinica, la maggioranza dei partecipanti la immaginava più intelligente e più competente di una fiduciosa. I ricercatori l'hanno chiamata «illusione del genio cinico».\footnote{Stavrova, O., e Ehlebracht, D. (2019). The cynical genius illusion: Exploring and debunking lay beliefs about cynicism and competence. \textit{Personality and Social Psychology Bulletin}, 45(2), 254-269.}

Chiamerò \index[cognitive]{bias del cinismo}\textit{bias del cinismo} il passo successivo, quello che dal sospetto verso gli altri arriva alla resa. Siccome nessuno è perfettamente obiettivo, tanto vale non credere a niente. Il suo effetto più grave è che svuota la responsabilità. Se tutto è manipolazione, nessuno risponde di nulla, nemmeno io.

\section{Testa, pancia e petto}

Platone\index[main]{Platone} aveva disegnato una mappa dell'anima che conviene riprendere come metafora, non come anatomia.\footnote{Platone. \textit{Repubblica}, Libro IV, 435a-441c.} In alto la \textit{ragione}, la testa. In basso l'\textit{appetito}, la pancia. In mezzo il \index[main]{thymos}\textbf{thymos}, l'«animo», la parte che si indigna davanti a un'ingiustizia, che prova orgoglio e vergogna. È il cavallo bianco del carro alato incontrato nel cap.~\ref{cap:p1c1}. Francis Fukuyama, riprendendo Platone, lo descrive come la parte dell'anima che desidera il riconoscimento della propria dignità.\footnote{Fukuyama, F. (1992). \textit{The End of History and the Last Man}. Free Press, New York.} Per Platone la ragione da sola è troppo debole, e le serve un alleato, un thymos ben educato che davanti alla tentazione non fa un calcolo ma dice «non lo faccio, perché è indegno di me».

Nel 1943 C. S. Lewis riprese questa mappa in \textit{L'abolizione dell'uomo}.\footnote{Lewis, C. S. (1943). \textit{The Abolition of Man}. Oxford University Press, Oxford.} Nei libri di testo del suo tempo, per rendere i ragazzi immuni alla propaganda, si insegnava a diffidare dei sentimenti, come se chi chiama «sublime» una cascata parlasse solo del proprio stato d'animo. Lewis obiettava che un sentimento educato è una risposta \textit{adeguata} a qualcosa di reale, e che insegnare a sfiduciarlo produce menti più fragili, non più critiche. Il risultato sono «uomini senza petto», con la testa e la pancia ma niente in mezzo. «Un cuore duro», scriveva, «non costituisce affatto una protezione infallibile contro una testa molle.» Il compito dell'educatore, concludeva, non è abbattere giungle ma irrigare deserti.

\figurasola{images/giungla_deserto.png}{La metafora di Lewis: il pericolo non è la giungla emotiva in sé, ma il deserto che lasciamo quando la tagliamo senza educarla.}{fig:giungla_deserto}

\section{Che cosa c'entrano i bias}

Lewis parlava di educazione morale, ma la ricerca successiva gli dà un appoggio inatteso. Le emozioni non sono soltanto un disturbo del ragionamento, ne fanno parte (cap.~\ref{cap:p1c2}). Le persone usano il proprio stato d'animo come un'informazione sul mondo, anche quando ha tutt'altra origine. In uno studio classico, chi veniva intervistato in una giornata di sole si dichiarava più soddisfatto della propria vita di chi veniva intervistato in una giornata di pioggia, a meno che l'intervistatore non menzionasse il tempo.\footnote{Schwarz, N., e Clore, G. L. (1983). Mood, misattribution, and judgments of well-being. \textit{Journal of Personality and Social Psychology}, 45(3), 513-523. L'effetto del meteo è stato poi trovato più piccolo in alcune repliche; il principio generale, l'uso dell'umore come informazione, resta ben documentato.} È una forma semplice di \index[cognitive]{euristica dell'affetto}euristica dell'affetto. E i pazienti studiati da Antonio Damasio, che per lesioni cerebrali avevano perso parte della risposta emotiva, non diventavano decisori più razionali. Faticavano a decidere.\footnote{Damasio, A. (1994). \textit{Descartes' Error: Emotion, Reason, and the Human Brain}. Putnam, New York. Alcuni aspetti della sua «ipotesi del marcatore somatico» sono discussi, ma il ruolo delle emozioni nella decisione è ampiamente accettato.}

Il problema, allora, non sono le emozioni ma le emozioni non riconosciute. Un'emozione che sappiamo nominare («sono indignato, e so perché») è un segnale che possiamo valutare. Una che abbiamo imparato a disprezzare ci guida lo stesso, senza che ce ne accorgiamo, ed è la condizione in cui i bias lavorano meglio. Non a caso i messaggi carichi di linguaggio morale ed emotivo tendono a diffondersi di più sui social network.\footnote{Brady, W. J., Wills, J. A., Jost, J. T., Tucker, J. A., e Van Bavel, J. J. (2017). Emotion shapes the diffusion of moralized content in social networks. \textit{PNAS}, 114(28), 7313-7318.} Chi non sa leggere la propria indignazione la lascia leggere, e usare, a qualcun altro.

\section{Firmare il proprio giudizio}

Che cosa si può fare, e chi può farlo? I divieti discussi nel cap.~\ref{cap:p2c29} possono avere senso, ma da soli ripetono l'errore descritto da Lewis, credere che basti togliere. Il thymos si allena con pratiche modeste, alla portata di ciascuno e soprattutto di chi educa.

La prima è formulare il proprio giudizio \textit{prima} di consultare la macchina o il gruppo, anche in una sola riga, per avere qualcosa con cui confrontare la risposta invece di riceverla e basta. La seconda è dare un nome a ciò che si sente davanti a una notizia o a un consenso troppo unanime, e chiedersi se la reazione è proporzionata al suo oggetto. La terza è distinguere il «non lo so» dal «non importa». Il primo è onestà, il secondo è il cinismo che si traveste da prudenza. A scuola e in famiglia, tutto questo passa meno dalle prediche che dalle storie lette e discusse insieme, dove si impara che cosa merita stupore, indignazione, fedeltà.

I limiti vanno detti con la stessa chiarezza. Un thymos allenato non rende immuni dai bias, e non va confuso con l'ostinazione. Rende soltanto più difficile che qualcuno, umano o algoritmo, decida al posto nostro che cosa dovremmo sentire e pensare.

\bigskip

C'è però un rovescio. Chi ha imparato a firmare un giudizio proprio e a difenderlo contro la macchina e contro il gruppo acquista una sicurezza preziosa e, insieme, pericolosa. Il coraggio di dire «questo lo penso io» è vicinissimo alla convinzione di vedere più chiaro degli altri. E proprio quella convinzione, come la ricerca sul punto cieco suggerisce da tempo, è il terreno in cui i bias si nascondono meglio.
